\documentclass[11pt]{article}
\usepackage[margin=1in]{geometry}
\usepackage{amsmath,amssymb}
\usepackage{enumitem}
\usepackage{longtable}
\usepackage{booktabs}
\usepackage{array}
\usepackage[hidelinks]{hyperref}

\title{Random Matrix Theory in Data Science and Statistics\\[4pt]\large Full Course Design}
\author{}
\date{}

\begin{document}
\maketitle

\section{Course Overview}

\begin{description}[leftmargin=2.2cm,style=nextline]
\item[Course name] Random Matrix Theory in Data Science and Statistics
\item[Level] PhD
\item[Format] Lecture course
\item[Duration] 14--15 weeks, meeting twice weekly, 75 minutes per session (28--30 sessions total; this design uses 30 sessions and can be compressed to 28 by shortening the review session and combining one pair of Module 2 sessions)
\item[Assessment philosophy] A small number of written homework assignments (four, one per module), continuous class participation, and a final project -- an in-class presentation together with a written report -- on a recent research paper, an open problem, or a topic of interest related to the course material. The final project is individual work (assumed, as this is typical for PhD-level project deliverables of this kind, unless the instructor specifies otherwise).
\item[Assumed background] Graduate-level, measure-theoretic probability (random variables, expectation, moments, modes of convergence, independence, basic concentration inequalities such as Markov/Chebyshev/Hoeffding); linear algebra through eigenvalues, eigenvectors, the singular value decomposition, and matrix norms; and general mathematical maturity with analysis-style proofs. More specialized machinery -- the combinatorics of non-crossing partitions, the Stieltjes transform and elementary complex analysis, and elements of operator theory needed for the non-commutative Khintchine inequality -- is introduced within the course itself and is not assumed.
\end{description}

\section{Course Schedule}

\begin{longtable}{p{1.4cm}p{1.4cm}p{9.5cm}}
\toprule
\textbf{Session} & \textbf{Module} & \textbf{Topic} \\
\midrule
\endhead
1 & --- & Introduction: course overview, motivation, review of prerequisites \\
2 & M1 & 1.1 Geometric methods of analysis with concentration inequalities (I) \\
3 & M1 & 1.1 Geometric methods of analysis with concentration inequalities (II) \\
4 & M1 & 1.2 Invariance and relation to random projections \\
5 & M1 & 1.3 Johnson--Lindenstrauss transform and dimensionality reduction \\
6 & M1 & 1.4 Compressed sensing \\
7 & M2 & 2.1 Moment method for eigenvalue limit theorems (I) -- \emph{Homework 1 due} \\
8 & M2 & 2.1 Moment method for eigenvalue limit theorems (II): combinatorics \\
9 & M2 & 2.2 Semicircle law for Wigner matrices \\
10 & M2 & 2.3 Marchenko--Pastur law for Wishart matrices (I) \\
11 & M2 & 2.3 Marchenko--Pastur law for Wishart matrices (II) \\
12 & M2 & 2.4 Elements of universality \\
13 & M2 & 2.5 Elements of free probability theory \\
14 & M2 & 2.6 Application: neural networks and random optimization landscapes \\
15 & M2 & 2.7 Application: covariance estimation \\
16 & M3 & 3.1 Resolvent method for eigenvalue limit theorems (I) -- \emph{Homework 2 due; project proposal due} \\
17 & M3 & 3.1 Resolvent method for eigenvalue limit theorems (II): eigenvectors \\
18 & M3 & 3.2 BBP phase transition in the performance of PCA \\
19 & M3 & 3.3 Computational challenges and statistical-computational gaps in PCA \\
20 & M3 & 3.4 Application: non-linear improvements to PCA \\
21 & M3 & 3.5 Application: community detection \\
22 & M4 & 4.1 General concentration inequalities (bounded differences, Lipschitz concentration) -- \emph{Homework 3 due} \\
23 & M4 & 4.2 Matrix Chernoff and Bernstein inequalities (I) \\
24 & M4 & 4.2 Matrix Chernoff and Bernstein inequalities (II): worked examples \\
25 & M4 & 4.3 Non-commutative Khintchine inequality and its extensions \\
26 & M4 & 4.4 Application: randomized numerical linear algebra \\
27 & --- & Review and synthesis -- \emph{Homework 4 due} \\
28 & --- & Final project presentations (I) \\
29 & --- & Final project presentations (II) \\
30 & --- & Final project presentations (III) -- \emph{final report due} \\
\bottomrule
\end{longtable}

%=================================================================
\section{Module 1: Gaussian Matrices and Dimensionality Reduction}
\emph{Goal: establish concentration-of-measure and random-projection techniques for analyzing Gaussian random matrices, and show how they yield dimension-reduction guarantees.}

\subsection{Geometric methods of analysis with concentration inequalities}
\textbf{Core concepts}
\begin{itemize}[nosep]
\item The concentration-of-measure phenomenon for high-dimensional Gaussian vectors
\item The Gaussian concentration inequality for Lipschitz functions
\item Sub-gaussian random variables and tail bounds
\item Covering numbers and $\varepsilon$-net arguments
\end{itemize}
\textbf{Learning objectives}
\begin{itemize}[nosep]
\item Explain the concentration-of-measure phenomenon and why high-dimensional Gaussian vectors concentrate near a sphere
\item Derive the Gaussian concentration inequality for Lipschitz functions
\item Apply concentration bounds to bound the extreme singular values of a Gaussian random matrix
\item Evaluate when a naive union bound is too weak and a covering/concentration argument is needed
\end{itemize}
\textbf{Difficulty note.} \emph{Why hard:} notational overload from indexing walks and pairings, combined with an abstraction jump from ``counting'' to ``which structures survive as $n\to\infty$'' resurfaces here in miniature through covering-number bookkeeping. \emph{Mitigation:} work a fully explicit $\varepsilon$-net construction on a low-dimensional sphere by hand before stating the general covering bound, and separate the ``fixed-vector concentration'' step from the ``union bound over the net'' step explicitly on the board.\\[4pt]
\textbf{Example questions}
\begin{enumerate}[nosep]
\item \emph{(Easier)} State the Gaussian concentration inequality for a Lipschitz function and use it to bound $P(|\lVert g\rVert - \sqrt n| > t)$ for a standard Gaussian vector $g \in \mathbb{R}^n$. \\
\emph{Rubric:} full credit requires citing $P(|f(g)-\mathbb{E}f(g)|>t) \le 2\exp(-t^2/2L^2)$ for $L$-Lipschitz $f$, identifying $f(x)=\lVert x\rVert$ as $1$-Lipschitz, and bounding $\mathbb{E}\lVert g\rVert$ to within $O(1/\sqrt n)$ of $\sqrt n$.
\item \emph{(Harder)} Use a covering-number argument together with Gaussian concentration to bound the largest singular value of an $n\times n$ Gaussian random matrix. \\
\emph{Rubric:} full credit requires an $\varepsilon$-net argument on the unit sphere, a union bound over the net combined with the concentration inequality, and a high-probability bound of the form $\sigma_{\max} \lesssim \sqrt n$.
\end{enumerate}

\subsection{Invariance and relation to random projections}
\textbf{Core concepts}
\begin{itemize}[nosep]
\item Rotational invariance of the Gaussian measure
\item Random projections as a model for dimensionality reduction
\item Distributional equivalence between random projections and Gaussian random matrices
\end{itemize}
\textbf{Learning objectives}
\begin{itemize}[nosep]
\item Explain why Gaussian matrices are rotationally invariant and what this buys analytically
\item Derive the distribution of the projection of a fixed vector onto a random subspace
\item Apply invariance arguments to simplify high-dimensional probability calculations
\end{itemize}
\textbf{Example questions}
\begin{enumerate}[nosep]
\item \emph{(Easier)} Show that if $g$ is a standard Gaussian vector in $\mathbb{R}^n$ and $U$ is any orthogonal matrix, then $Ug$ has the same distribution as $g$. \\
\emph{Rubric:} full credit for correctly using the invariance of the Gaussian density (or characteristic function) under orthogonal transformations.
\item \emph{(Harder)} Derive the distribution of the squared length of the projection of a fixed unit vector onto a uniformly random $k$-dimensional subspace of $\mathbb{R}^n$. \\
\emph{Rubric:} full credit requires reducing to a Beta (or scaled ratio-of-chi-squares) distribution via rotational invariance, with correctly identified parameters in terms of $k$ and $n$.
\end{enumerate}

\subsection{Johnson--Lindenstrauss transform and dimensionality reduction}
\textbf{Core concepts}
\begin{itemize}[nosep]
\item Statement of the Johnson--Lindenstrauss (JL) lemma
\item Random projection as a realization of a JL transform
\item Trade-off between target dimension, distortion, and failure probability
\end{itemize}
\textbf{Learning objectives}
\begin{itemize}[nosep]
\item Explain the statement and significance of the JL lemma
\item Derive the JL lemma using Gaussian concentration and a union bound
\item Apply the JL transform to determine the projection dimension needed for a given distortion and failure tolerance
\item Evaluate the optimality of JL-type dimension bounds (matching lower bounds, sparse JL variants)
\end{itemize}
\textbf{Example questions}
\begin{enumerate}[nosep]
\item \emph{(Easier)} State the Johnson--Lindenstrauss lemma and explain the role of each parameter ($n$, $k$, $\varepsilon$, $\delta$). \\
\emph{Rubric:} full credit for a precise statement and a correct explanation of each parameter's role.
\item \emph{(Harder)} Derive the minimum projection dimension $k$ needed to preserve pairwise distances among $m$ points to within distortion $\varepsilon$ with failure probability at most $\delta$. \\
\emph{Rubric:} full credit requires combining a per-pair concentration bound with a union bound over $O(m^2)$ pairs, yielding $k = O(\varepsilon^{-2}\log(m/\delta))$.
\end{enumerate}

\subsection{Compressed sensing}
\textbf{Core concepts}
\begin{itemize}[nosep]
\item Sparse signal recovery
\item The restricted isometry property (RIP)
\item Connection between RIP and Gaussian/subgaussian random matrices
\item $\ell_1$-minimization recovery guarantees
\end{itemize}
\textbf{Learning objectives}
\begin{itemize}[nosep]
\item Explain the compressed sensing problem and the role of sparsity
\item Derive that a suitably scaled Gaussian matrix satisfies the RIP with high probability using covering/concentration arguments
\item Apply RIP-based recovery guarantees to bound the number of measurements needed for exact recovery
\item Evaluate why $\ell_0$-minimization is intractable and how $\ell_1$-relaxation circumvents this
\end{itemize}
\textbf{Example questions}
\begin{enumerate}[nosep]
\item \emph{(Easier)} Define the restricted isometry property (RIP) of order $s$ and explain its role in sparse recovery. \\
\emph{Rubric:} full credit for the correct RIP definition and a correct explanation of its connection to stable recovery guarantees.
\item \emph{(Harder)} Sketch a proof (via covering plus union bound) that a suitably scaled Gaussian matrix satisfies the RIP of order $s$ with high probability, and state the resulting bound on the number of measurements. \\
\emph{Rubric:} full credit requires correctly bounding the number of $s$-sparse subspaces, applying a per-subspace concentration bound, and arriving at $m = O(s\log(n/s))$.
\end{enumerate}

%=================================================================
\section{Module 2: Classical Theory of i.i.d. Random Matrices}
\emph{Goal: develop the classical spectral limit theorems (semicircle, Marchenko--Pastur) for i.i.d. random matrices via moment and combinatorial methods, and connect them to free probability and data-science applications.}

\subsection{Moment method for eigenvalue limit theorems}
\textbf{Core concepts}
\begin{itemize}[nosep]
\item The method of moments for weak convergence of empirical spectral distributions
\item Traces of matrix powers as sums over closed walks
\item Combinatorics of pairings and non-crossing partitions
\end{itemize}
\textbf{Learning objectives}
\begin{itemize}[nosep]
\item Explain why matching moments implies convergence of the empirical spectral distribution, under suitable conditions
\item Derive the moment formula for $\mathbb{E}[\mathrm{tr}(M^n)]$ in terms of combinatorial structures (walks/pairings)
\item Apply the moment method to compute low-order moments of a random matrix ensemble
\item Evaluate which combinatorial structures survive in the large-$n$ limit and which vanish
\end{itemize}
\textbf{Difficulty note.} \emph{Why hard:} notational overload from indexing walks/pairings and tracking vertex identifications, plus an abstraction jump from ``counting'' to ``which structures survive as $n \to \infty$.'' \emph{Mitigation:} work a fully explicit small example (e.g., compute $\mathbb{E}[\mathrm{tr}(M^4)]$ by hand for a $2\times2$ or $3\times3$ Wigner matrix) before stating the general combinatorial formula, and use a diagram-drawing exercise to visualize pairings.\\[4pt]
\textbf{Example questions}
\begin{enumerate}[nosep]
\item \emph{(Easier)} Compute $\mathbb{E}[\mathrm{tr}(M^n)]$ for small $n$ (e.g., $n=2,4$) for a Wigner matrix by directly expanding the trace as a sum over closed walks. \\
\emph{Rubric:} full credit for correctly expanding the sum and identifying which terms survive in expectation given independence and mean-zero entries.
\item \emph{(Harder)} Explain why, in the large-matrix limit, only non-crossing pairings of a closed walk contribute to the leading-order moment, briefly justifying this via the scaling of the number of pairings against matrix dimension. \\
\emph{Rubric:} full credit requires the counting argument (non-crossing vs.\ total pairings) and the matching power-of-$n$ scaling.
\end{enumerate}

\subsection{Semicircle law for Wigner matrices}
\textbf{Core concepts}
\begin{itemize}[nosep]
\item Wigner matrix ensembles
\item The semicircle distribution and its density
\item Non-crossing pairings and Catalan numbers
\end{itemize}
\textbf{Learning objectives}
\begin{itemize}[nosep]
\item Explain the statement of Wigner's semicircle law
\item Derive the moments of the semicircle distribution via Catalan-number counting of non-crossing pairings
\item Apply the moment method to prove convergence of the empirical spectral distribution to the semicircle law
\item Evaluate the universality of the semicircle law across different entry distributions
\end{itemize}
\textbf{Example questions}
\begin{enumerate}[nosep]
\item \emph{(Easier)} State Wigner's semicircle law and give the density of the semicircle distribution. \\
\emph{Rubric:} full credit for the correct statement and correct density formula.
\item \emph{(Harder)} Show that the even moments of the semicircle distribution are given by the Catalan numbers, and the odd moments vanish; explain the combinatorial reason. \\
\emph{Rubric:} full credit requires the correct Catalan-number formula and its connection to counting non-crossing pairings.
\end{enumerate}

\subsection{Marchenko--Pastur law for Wishart matrices}
\textbf{Core concepts}
\begin{itemize}[nosep]
\item Sample covariance (Wishart) matrices
\item The high-dimensional aspect-ratio regime
\item The Marchenko--Pastur distribution and its support
\end{itemize}
\textbf{Learning objectives}
\begin{itemize}[nosep]
\item Explain the setup of Wishart matrices and the high-dimensional aspect-ratio regime
\item Derive the moments of the Marchenko--Pastur distribution
\item Apply the Marchenko--Pastur law to describe the spectrum of a sample covariance matrix in high dimensions
\item Evaluate the implications of the Marchenko--Pastur law for the bias of naive sample-eigenvalue estimates
\end{itemize}
\textbf{Example questions}
\begin{enumerate}[nosep]
\item \emph{(Easier)} State the Marchenko--Pastur law and describe how its support depends on the aspect ratio. \\
\emph{Rubric:} full credit for the correct density/support formula in terms of the aspect-ratio parameter.
\item \emph{(Harder)} Derive the first two moments of the Marchenko--Pastur distribution via the moment method and verify they match a direct computation of $\mathbb{E}[\mathrm{tr}(W)]$ and $\mathbb{E}[\mathrm{tr}(W^2)]$ for a Wishart matrix $W$. \\
\emph{Rubric:} full credit requires a correct moment computation by both routes and a consistency check between them.
\end{enumerate}

\subsection{Elements of universality}
\textbf{Core concepts}
\begin{itemize}[nosep]
\item The universality principle for random matrix spectra
\item Four-moment theorems (conceptual level)
\item Comparison / Lindeberg-swapping arguments (conceptual level)
\end{itemize}
\textbf{Learning objectives}
\begin{itemize}[nosep]
\item Explain the universality principle and precisely what it does and does not claim
\item Compare proof strategies (moment method vs.\ comparison/swapping arguments) for establishing universality
\item Evaluate the evidence for or against universality in a given, possibly non-Gaussian, matrix ensemble
\end{itemize}
\textbf{Example questions}
\begin{enumerate}[nosep]
\item \emph{(Easier)} Explain what the universality principle claims about the applicability of the semicircle law beyond Gaussian entries. \\
\emph{Rubric:} full credit for correctly stating that the limiting spectral distribution depends only on the first two moments of the entries, not the full entry distribution.
\item \emph{(Harder)} Compare the moment-method proof of universality with a comparison/Lindeberg-swapping argument, identifying one advantage of each. \\
\emph{Rubric:} full credit requires a substantive, accurate comparison (e.g., the moment method gives explicit control but needs all moments to exist; swapping arguments need fewer moment assumptions but are less explicit).
\end{enumerate}

\subsection{Elements of free probability theory}
\textbf{Core concepts}
\begin{itemize}[nosep]
\item Free independence vs.\ classical independence
\item Free convolution
\item The $R$-transform and its role in computing spectra of matrix sums
\end{itemize}
\textbf{Learning objectives}
\begin{itemize}[nosep]
\item Explain the notion of freeness and how it differs from classical independence
\item Derive the free convolution of two simple spectral distributions using the $R$-transform
\item Apply free probability to predict the limiting spectrum of a sum of independent random matrices
\item Evaluate when free-probability heuristics correctly predict finite-$n$ behavior
\end{itemize}
\textbf{Difficulty note.} \emph{Why hard:} freeness is a genuinely new axiomatic framework that looks superficially like independence but is not, and $R$-transform manipulations can feel like formal algebra with no intuition. \emph{Mitigation:} repeatedly contrast a free-probability computation with its classical-probability analogue side by side, and use the semicircle-plus-semicircle example as a running worked case before stating the general $R$-transform machinery.\\[4pt]
\textbf{Example questions}
\begin{enumerate}[nosep]
\item \emph{(Easier)} Define free independence and explain how it differs from classical independence for two random variables/matrices. \\
\emph{Rubric:} full credit for correctly stating the mixed-moment condition defining freeness and contrasting it with classical independence.
\item \emph{(Harder)} Use the $R$-transform to compute the limiting spectral distribution of the sum of two independent semicircle-distributed random matrices. \\
\emph{Rubric:} full credit requires correctly using additivity of the $R$-transform under free convolution and recovering a (rescaled) semicircle law.
\end{enumerate}

\subsection{Application: neural networks and random optimization landscapes}
\textbf{Core concepts}
\begin{itemize}[nosep]
\item Random-matrix models of neural-network loss Hessians
\item Spectral properties and optimization-landscape geometry (saddle points vs.\ minima)
\end{itemize}
\textbf{Learning objectives}
\begin{itemize}[nosep]
\item Explain how random matrix spectra are used to model the Hessian of a high-dimensional loss landscape
\item Apply semicircle/Marchenko--Pastur-type results to characterize the proportion of negative eigenvalues (saddle directions) in a toy model
\item Evaluate the strengths and limitations of random-matrix models of real neural-network landscapes
\end{itemize}
\textbf{Example questions}
\begin{enumerate}[nosep]
\item \emph{(Easier)} Explain how a random-matrix model of a loss Hessian can be used to estimate the fraction of negative eigenvalues at a given energy level. \\
\emph{Rubric:} full credit for a correct qualitative explanation connecting the spectral density to the fraction of negative eigenvalues.
\item \emph{(Harder)} For a toy Hessian model built from a Wigner matrix plus a shift, compute the fraction of negative eigenvalues as a function of the shift using the semicircle law. \\
\emph{Rubric:} full credit for correctly integrating the shifted semicircle density over the negative axis.
\end{enumerate}

\subsection{Application: covariance estimation}
\textbf{Core concepts}
\begin{itemize}[nosep]
\item Sample covariance estimation in high dimensions
\item Eigenvalue shrinkage
\item Consequences of the Marchenko--Pastur law for estimator design
\end{itemize}
\textbf{Learning objectives}
\begin{itemize}[nosep]
\item Explain why the sample covariance matrix is an inconsistent estimator of eigenvalues in high dimensions
\item Apply Marchenko--Pastur-based corrections (e.g., eigenvalue shrinkage) to design an improved covariance estimator
\item Evaluate the estimator's performance against the naive sample covariance in a high-dimensional regime
\end{itemize}
\textbf{Example questions}
\begin{enumerate}[nosep]
\item \emph{(Easier)} Explain why the largest eigenvalue of a sample covariance matrix is a biased estimator of the largest population eigenvalue in the fixed-aspect-ratio, high-dimensional regime. \\
\emph{Rubric:} full credit for connecting this to Marchenko--Pastur-induced spreading of the sample spectrum.
\item \emph{(Harder)} Propose and justify a simple eigenvalue-shrinkage correction to the sample covariance matrix using the Marchenko--Pastur law, for an isotropic population covariance. \\
\emph{Rubric:} full credit requires a correct derivation connecting the MP-predicted spectrum to a shrinkage rule and a justification that it removes the bias.
\end{enumerate}

%=================================================================
\section{Module 3: Spiked Matrix Models and Principal Component Analysis}
\emph{Goal: analyze low-rank (``spiked'') perturbations of random matrices via the resolvent method to characterize the BBP phase transition and its implications for PCA, including statistical-computational tradeoffs.}

\subsection{Resolvent method for eigenvalue limit theorems and eigenvectors}
\textbf{Core concepts}
\begin{itemize}[nosep]
\item The resolvent (Stieltjes transform) of a random matrix
\item Self-consistent (fixed-point) equations for the resolvent
\item Recovering the spectral density and eigenvector information from the resolvent
\end{itemize}
\textbf{Learning objectives}
\begin{itemize}[nosep]
\item Explain the definition and role of the resolvent/Stieltjes transform
\item Derive a self-consistent equation for the resolvent of a simple random matrix ensemble
\item Apply the resolvent method to recover the semicircle or Marchenko--Pastur law as an alternative to the moment method
\item Evaluate the advantages of the resolvent method over the moment method (e.g., access to eigenvector information)
\end{itemize}
\textbf{Difficulty note.} \emph{Why hard:} the complex-analysis background (analytic continuation) needed is often not assumed, and there is notational overload from tracking several related transforms (Stieltjes, Cauchy, $R$-transform). \emph{Mitigation:} dedicate the first part of the session to a self-contained mini-primer on exactly the complex-analysis facts needed, and immediately connect the resolvent back to the already-familiar semicircle law from Module 2 as a consistency check.\\[4pt]
\textbf{Example questions}
\begin{enumerate}[nosep]
\item \emph{(Easier)} Define the resolvent (Stieltjes transform) of a random matrix and state its relationship to the empirical spectral distribution. \\
\emph{Rubric:} full credit for the correct definition and correct inversion relationship.
\item \emph{(Harder)} Derive the self-consistent equation for the resolvent of a Wigner matrix and use it to recover the semicircle law. \\
\emph{Rubric:} full credit requires a correct derivation of the fixed-point equation and correct algebraic recovery of the semicircle density.
\end{enumerate}

\subsection{BBP phase transition in the performance of PCA}
\textbf{Core concepts}
\begin{itemize}[nosep]
\item Spiked covariance/Wigner models
\item The critical signal-strength threshold
\item Eigenvalue outlier and eigenvector-alignment behavior above/below threshold
\end{itemize}
\textbf{Learning objectives}
\begin{itemize}[nosep]
\item Explain the spiked-model setup and what ``signal strength'' means in this context
\item Derive the critical threshold at which an outlier eigenvalue emerges from the bulk
\item Apply the BBP result to predict whether PCA will detect a given signal at a given aspect ratio
\item Evaluate the practical implications of the BBP transition for the reliability of PCA in high dimensions
\end{itemize}
\textbf{Difficulty note.} \emph{Why hard:} the exact threshold formula and the ``flat region'' below threshold look unmotivated on first exposure, and it is easy to conflate ``an outlier eigenvalue exists'' with ``the eigenvector recovers the signal direction.'' \emph{Mitigation:} build the threshold up from a numerical exploration first (vary signal strength, plot the outlier location) before deriving it analytically, and explicitly separate the eigenvalue question from the eigenvector-alignment question as two distinct claims.\\[4pt]
\textbf{Example questions}
\begin{enumerate}[nosep]
\item \emph{(Easier)} State the BBP phase transition result for a rank-one spiked Wigner model, including the critical signal-strength threshold. \\
\emph{Rubric:} full credit for the correct threshold formula and a correct description of above/below-threshold behavior.
\item \emph{(Harder)} Derive the location of the outlier eigenvalue above the BBP threshold using the resolvent/self-consistent-equation approach. \\
\emph{Rubric:} full credit requires correctly setting up and solving the fixed-point equation for the outlier location as a function of signal strength.
\end{enumerate}

\subsection{Computational challenges and statistical-computational gaps in PCA}
\textbf{Core concepts}
\begin{itemize}[nosep]
\item Statistical (information-theoretic) vs.\ computational detection thresholds
\item Sparse PCA as a canonical example of a statistical-computational gap
\item Reduction-based and low-degree-polynomial evidence for hardness (conceptual level)
\end{itemize}
\textbf{Learning objectives}
\begin{itemize}[nosep]
\item Explain the distinction between statistical and computational detectability
\item Compare the information-theoretic threshold to the best known polynomial-time threshold for sparse PCA
\item Evaluate the evidence (average-case reductions, low-degree heuristics) supporting a conjectured statistical-computational gap
\end{itemize}
\textbf{Difficulty note.} \emph{Why hard:} the distinction between worst-case hardness, average-case hardness, and ``no known polynomial-time algorithm achieves the statistical threshold'' is conceptually slippery, and students often collapse these into one notion. \emph{Mitigation:} present a concrete side-by-side table (statistical threshold vs.\ best-known-polynomial-time threshold vs.\ worst-case NP-hardness) for the sparse-PCA example specifically, so the three notions stay visibly distinct.\\[4pt]
\textbf{Example questions}
\begin{enumerate}[nosep]
\item \emph{(Easier)} Explain the difference between the statistical detection threshold and the best-known polynomial-time detection threshold for sparse PCA. \\
\emph{Rubric:} full credit for a correct, precise distinction between the two thresholds.
\item \emph{(Harder)} Discuss one piece of evidence (e.g., a reduction from planted clique, or a low-degree-polynomial lower bound) supporting the conjecture that no polynomial-time algorithm closes the sparse-PCA statistical-computational gap. \\
\emph{Rubric:} full credit requires a substantively correct description of the evidence and an acknowledgment of its limits (evidence, not proof).
\end{enumerate}

\subsection{Application: non-linear improvements to PCA}
\textbf{Core concepts}
\begin{itemize}[nosep]
\item Limitations of linear PCA under non-Gaussian or structured signal models
\item Non-linear/pre-processing techniques (e.g., entrywise transforms) that move signal strength closer to the statistical threshold
\end{itemize}
\textbf{Learning objectives}
\begin{itemize}[nosep]
\item Explain why a well-chosen non-linear transform can improve estimation over linear PCA
\item Apply a non-linear pre-processing method to a spiked model and predict its effect on the effective signal strength
\item Evaluate the trade-offs (robustness, assumptions required) of non-linear PCA variants versus standard PCA
\end{itemize}
\textbf{Example questions}
\begin{enumerate}[nosep]
\item \emph{(Easier)} Explain why applying a non-linear entrywise transform before running PCA can sometimes improve signal detection compared to standard linear PCA. \\
\emph{Rubric:} full credit for correctly explaining the transform's effect on effective signal strength relative to the BBP threshold.
\item \emph{(Harder)} For a simple, specified non-linear transform, compute or estimate its effect on the effective spike strength in a toy spiked model. \\
\emph{Rubric:} full credit requires a correct (possibly approximate) computation connecting the transform to a change in effective signal-to-noise ratio.
\end{enumerate}

\subsection{Application: community detection}
\textbf{Core concepts}
\begin{itemize}[nosep]
\item The stochastic block model
\item Adjacency/Laplacian spectra and spectral clustering algorithms
\item Connection to spiked random matrix models
\end{itemize}
\textbf{Learning objectives}
\begin{itemize}[nosep]
\item Explain how community detection in the stochastic block model can be cast as a spiked-matrix recovery problem
\item Apply BBP-type thresholds to determine the detectability regime for a given block model
\item Evaluate the performance of spectral clustering against the information-theoretic detection threshold
\end{itemize}
\textbf{Example questions}
\begin{enumerate}[nosep]
\item \emph{(Easier)} Explain how the stochastic block model with two communities can be related to a spiked random matrix model. \\
\emph{Rubric:} full credit for correctly mapping block-model parameters (within/between-community edge probabilities) to a spike strength.
\item \emph{(Harder)} Using the BBP threshold, derive the condition on the block-model parameters under which spectral clustering can detect the community structure. \\
\emph{Rubric:} full credit requires correctly translating the BBP threshold condition into block-model parameter language.
\end{enumerate}

%=================================================================
\section{Module 4: Matrix Concentration Inequalities}
\emph{Goal: derive general-purpose non-asymptotic concentration inequalities for random matrices (matrix Chernoff/Bernstein, non-commutative Khintchine) and apply them to randomized numerical linear algebra.}

\subsection{General concentration inequalities and applications to random matrix norms}
\textbf{Core concepts}
\begin{itemize}[nosep]
\item McDiarmid's bounded-differences inequality
\item Lipschitz concentration for general (non-Gaussian) measures
\item Sub-gaussian and sub-exponential tail behavior
\end{itemize}
\textbf{Learning objectives}
\begin{itemize}[nosep]
\item Explain the bounded-differences condition and what it captures about a function's sensitivity
\item Derive McDiarmid's inequality, or apply its proof strategy, to a specific function of independent inputs
\item Apply Lipschitz concentration to bound fluctuations of a random matrix functional (e.g., an extreme singular value)
\end{itemize}
\textbf{Example questions}
\begin{enumerate}[nosep]
\item \emph{(Easier)} State McDiarmid's bounded-differences inequality and identify the bounded-differences condition for a simple function of independent inputs. \\
\emph{Rubric:} full credit for the correct statement and correct verification of the bounded-differences condition for the given example.
\item \emph{(Harder)} Apply Lipschitz concentration to bound the fluctuation of the largest singular value of a random matrix with independent, sub-gaussian (not necessarily Gaussian) entries. \\
\emph{Rubric:} full credit requires correctly verifying the Lipschitz condition for the singular-value functional and applying the concentration bound.
\end{enumerate}

\subsection{General-purpose bounds on expected random matrix norms}
\textbf{Core concepts}
\begin{itemize}[nosep]
\item Matrix versions of the Chernoff and Bernstein inequalities
\item The matrix variance proxy and matrix moment generating function
\item Intrinsic (effective) dimension in matrix concentration bounds
\end{itemize}
\textbf{Learning objectives}
\begin{itemize}[nosep]
\item Explain how scalar Chernoff/Bernstein inequalities generalize to the matrix setting
\item Derive a matrix Chernoff or Bernstein bound for a sum of independent random matrices using the matrix moment generating function
\item Apply matrix Bernstein to bound the norm of a random matrix arising from sampling or sparsification
\end{itemize}
\textbf{Difficulty note.} \emph{Why hard:} multiple competing definitions of the matrix moment generating function and ``variance proxy'' create notational overload, and the role of intrinsic/effective dimension is easy to miss. \emph{Mitigation:} derive the scalar Bernstein inequality first as a warm-up, then show step by step exactly which scalar quantity is replaced by which matrix quantity, explicitly flagging intrinsic dimension as the one genuinely new ingredient.\\[4pt]
\textbf{Example questions}
\begin{enumerate}[nosep]
\item \emph{(Easier)} State the matrix Bernstein inequality and identify the role of each term (variance proxy, dimension factor). \\
\emph{Rubric:} full credit for a correct statement and correct explanation of each term.
\item \emph{(Harder)} Apply matrix Bernstein to bound the norm of a sum of independent random rank-one matrices arising from sampling rows of a fixed matrix. \\
\emph{Rubric:} full credit requires correctly computing the matrix variance term and applying the inequality to obtain a high-probability norm bound.
\end{enumerate}

\subsection{Non-commutative Khintchine inequality and its extensions}
\textbf{Core concepts}
\begin{itemize}[nosep]
\item Statement of the non-commutative Khintchine inequality
\item Its connection to matrix concentration bounds
\item Recent sharpenings and extensions (conceptual level)
\end{itemize}
\textbf{Learning objectives}
\begin{itemize}[nosep]
\item Explain the statement of the non-commutative Khintchine inequality and its role as a sharper alternative to matrix Bernstein
\item Compare the non-commutative Khintchine bound to matrix Chernoff/Bernstein bounds in a specific example
\item Evaluate the significance of recent extensions of the inequality for random matrix theory
\end{itemize}
\textbf{Difficulty note.} \emph{Why hard:} the inequality relies on operator-space/Schatten-norm background most students have not seen, so it initially looks like an unmotivated formal statement. \emph{Mitigation:} state the classical (commutative) Khintchine inequality first as an anchor, then introduce the non-commutative version as ``the same statement with absolute values replaced by Schatten norms,'' deferring the technical operator-space justification to a reference.\\[4pt]
\textbf{Example questions}
\begin{enumerate}[nosep]
\item \emph{(Easier)} State the non-commutative Khintchine inequality and explain how it generalizes the classical (scalar) Khintchine inequality. \\
\emph{Rubric:} full credit for a correct statement and correct explanation of the generalization.
\item \emph{(Harder)} Compare the bound given by the non-commutative Khintchine inequality to the matrix Bernstein bound for a specific sum of random matrices, identifying which is sharper. \\
\emph{Rubric:} full credit requires a correct comparison with correct identification of the sharper bound and why.
\end{enumerate}

\subsection{Application: randomized numerical linear algebra}
\textbf{Core concepts}
\begin{itemize}[nosep]
\item Randomized sketching for matrix multiplication and regression
\item Approximation guarantees derived from matrix concentration
\item Trade-off between sketch size, approximation error, and failure probability
\end{itemize}
\textbf{Learning objectives}
\begin{itemize}[nosep]
\item Explain how random sketching reduces the cost of a numerical linear algebra task
\item Apply matrix concentration inequalities to derive an approximation guarantee for a sketched matrix computation
\item Design a randomized algorithm (e.g., for approximate matrix multiplication) with a specified error tolerance, using matrix concentration bounds to select the sketch size
\end{itemize}
\textbf{Example questions}
\begin{enumerate}[nosep]
\item \emph{(Easier)} Explain how a random sketching matrix can be used to approximately solve a large least-squares regression problem, and what guarantee is desired. \\
\emph{Rubric:} full credit for a correct explanation of the sketch-and-solve approach and the approximation guarantee sought.
\item \emph{(Harder)} Using a matrix concentration inequality, derive the sketch size needed to guarantee a $(1\pm\varepsilon)$ relative-error approximate matrix multiplication with high probability. \\
\emph{Rubric:} full credit requires correctly applying a matrix concentration bound to derive the sketch-size/error/failure-probability trade-off.
\end{enumerate}

%=================================================================
\section{Prerequisite Map}

\begin{longtable}{p{4.3cm}p{5.3cm}p{4.7cm}}
\toprule
\textbf{Topic} & \textbf{Prerequisite(s)} & \textbf{Where taught} \\
\midrule
\endhead
1.1 Geometric methods / concentration & Basic probability (tail bounds, sub-gaussian variables); linear algebra (norms, singular values) & Assumed background \\
1.2 Invariance / random projections & 1.1; rotational invariance of the Gaussian measure & 1.1 (this course); assumed background (linear algebra) \\
1.3 JL transform & 1.1; 1.2 & This course \\
1.4 Compressed sensing & 1.1; 1.3 (RIP builds on concentration and covering arguments) & This course \\
2.1 Moment method & Linear algebra (trace, eigenvalues); basic combinatorics & Assumed background (linear algebra); combinatorics introduced in 2.1 \\
2.2 Semicircle law & 2.1 & This course \\
2.3 Marchenko--Pastur law & 2.1 & This course \\
2.4 Universality & 2.2; 2.3 & This course \\
2.5 Free probability & 2.1; 2.2 (moment/combinatorial intuition) & This course \\
2.6 App: neural networks & 2.2; 2.3 (spectral results) & This course \\
2.7 App: covariance estimation & 2.3 & This course \\
3.1 Resolvent method & 2.1 (as contrast); basic complex analysis & This course (complex-analysis basics introduced here); assumed background (analysis maturity) \\
3.2 BBP transition & 3.1; 2.2 / 2.3 (bulk spectral laws) & This course \\
3.3 Statistical-computational gaps & 3.2 & This course \\
3.4 App: non-linear PCA & 3.2; 3.3 & This course \\
3.5 App: community detection & 3.2 & This course \\
4.1 Bounded-differences / Lipschitz concentration & 1.1 (concentration foundations) & This course (built on 1.1) \\
4.2 Matrix Chernoff / Bernstein & 4.1; linear algebra (matrix norms, trace) & This course (built on 4.1); assumed background (linear algebra) \\
4.3 Non-commutative Khintchine & 4.2 & This course \\
4.4 App: randomized numerical linear algebra & 4.2; 4.3 & This course \\
\bottomrule
\end{longtable}

%=================================================================
\section{Assessment Plan}

\begin{longtable}{p{3.5cm}p{3.5cm}p{2.3cm}p{5.0cm}}
\toprule
\textbf{Module / Milestone} & \textbf{Assessment} & \textbf{Timing} & \textbf{Rationale} \\
\midrule
\endhead
Module 1 & Homework 1 (Gaussian concentration, invariance, JL, compressed sensing) & Due session 7 & Tests students' ability to carry out concentration-based derivations independently before moving to the more demanding combinatorial machinery of Module 2 \\
Module 2 & Homework 2 (moment method, semicircle law, Marchenko--Pastur, universality, free probability) & Due session 16 & Tests mastery of the combinatorial/moment machinery that underlies the resolvent-based methods used in Module 3 \\
Module 2 & Final project proposal & Due session 16 & Ensures students select a paper, open problem, or topic early enough to pursue a genuine deep dive alongside coursework \\
Module 3 & Homework 3 (resolvent method, BBP transition, statistical-computational gaps) & Due session 22 & Tests the resolvent technique and phase-transition reasoning central to spiked-matrix models \\
Module 4 & Homework 4 (matrix concentration, Chernoff/Bernstein, non-commutative Khintchine) & Due session 27 & Tests the general-purpose, non-asymptotic concentration toolkit and its use in an application \\
All modules & Class participation & Continuous, every session & Given the lecture format, participation credit (questions, engagement with board work and derivations) incentivizes active rather than passive attendance \\
Capstone & Final project: in-class presentation & Sessions 28--30 & Gives students practice presenting technical research material and exposes the class to material beyond the syllabus \\
Capstone & Final project: written report & Due session 30 & Assesses depth of independent engagement with a research paper, open problem, or related topic beyond what is covered in lecture \\
\bottomrule
\end{longtable}

\end{document}